\section{Design}
\label{sec:design}

\Cref{fig:arch} shows the parts. Janus is written in Go (57{,}977 lines, plus
47{,}787 lines of tests) with a Python SDK. % source: find pkg cmd -name '*.go' | xargs cat | wc -l, test and non-test files separately (2026-09-29)
The wire vocabulary is Protocol Buffers and is the same grammar the log records,
so a message and the evidence of it are one type.

\begin{figure}[t]
\centering
\begin{tikzpicture}[font=\scriptsize,
  box/.style={draw,rounded corners=2pt,align=center,minimum height=0.7cm,inner sep=2.5pt,fill=white},
  arr/.style={-{Latex[length=1.6mm]},thin}]
\node[box,minimum width=1.55cm] (agent) at (0,1.1) {LLM agent\\(SDK)};
\node[box,minimum width=1.55cm] (mcp) at (0,-0.1) {MCP\\proxy};
\node[box,minimum width=2.3cm,minimum height=1.9cm] (orchd) at (2.95,0.5)
  {\textbf{janus-orchd}\\[1pt]saga fold\\gates\\outbox\\registry};
\node[box,minimum width=1.55cm] (val) at (5.9,1.1) {validator /\\approver};
\node[box,minimum width=1.55cm] (tool) at (5.9,-0.1) {tool /\\effect};
\node[box,minimum width=6.9cm] (log) at (2.95,-1.55)
  {\textbf{evidence log}: CBOR envelopes, BLAKE3 chain, segments signed with\\Ed25519 over a Merkle root (RFC 6962 tree, BLAKE3)};
\node[box,minimum width=2.2cm] (proj) at (1.2,-2.75) {projections\\(rebuildable)};
\node[box,minimum width=2.2cm] (verify) at (4.7,-2.75) {janus-verify\\(auditor, offline)};
\draw[arr] (agent.east) -- node[above,pos=0.45]{propose} (agent.east-|orchd.west);
\draw[arr] (mcp.east) -- node[below,pos=0.45]{intercept} (mcp.east-|orchd.west);
\draw[arr] ([yshift=4pt]orchd.east|-val) -- node[above]{question} ([yshift=4pt]val.west);
\draw[arr] ([yshift=-4pt]val.west) -- node[below]{answer} ([yshift=-4pt]orchd.east|-val);
\draw[arr] (orchd.east|-tool) -- node[above]{2. release} (tool.west);
\draw[arr] (orchd.south) -- node[right]{1. append} (orchd.south|-log.north);
\draw[arr] (log.south-|proj) -- (proj.north);
\draw[arr] (log.south-|verify) -- (verify.north);
\end{tikzpicture}
\caption{Janus. A decision is appended (1) before the effect it permits is
released; answers from outside arrive as records, not calls. The log is the only
source of truth: projections are rebuilt from it, and an auditor needs only the
log and a public key.}
\label{fig:arch}
\end{figure}

\subsection{The evidence log}
\label{sec:design-log}
Each record is a CBOR envelope (kind, saga, step, participant, timestamps, the
hash of its payload) around a Protocol Buffers payload. Records are chained: a
record's chain hash is BLAKE3 over the previous chain hash, its payload hash and
its header. Records are grouped into segments, and a segment's footer carries a
Merkle root over its records and one Ed25519 signature binding that root to the
segment header, the record count and the sequence range. The tree has the
structure of RFC~6962~\S2.1 (since obsoleted by RFC~9162)~\cite{rfc6962,rfc9162},
instantiated with BLAKE3 rather than SHA-256, and is built per segment so that
signing is off the per-event path. A single event is proven by an inclusion
proof against its segment's signed root, and the verifier checks that the event a
bundle names is the record its proof covers. Rotations of the writer's key are
records in the log, each declared in a segment its predecessor signed, so an
auditor is handed one root key rather than a key list. The chain and the tree
are the tamper-evident logging literature's data
structures~\cite{schneierkelsey,crosbywallach,haberstornetta}; its defence against
a compromised logger is not implemented (\cref{sec:related,sec:model-out}).

A single writer appends, with group commit: concurrent callers are batched and
share one durability barrier. Everything else (the relational projections the
console and the gates query, the saga index, the registry view) is rebuilt from
the log and is never the source of truth; a projection answers only as far as
it has folded and says so.

\subsection{The saga engine}
A saga's state is a fold over its records, a pure function with no clock and no
I/O. The coordinator that drives it appends a record and only then acts on the
state that record implies, so a coordinator killed at any point resumes from the
log and reaches the same state. Compensation runs in reverse topological order
of the plan's dependencies (I7), sub-sagas either cascade to their parent's
commit or commit autonomously, and a frontier gate refuses a commit over an
unsealed conflicting claim by another saga on the same resource (I6).

\emph{Recorded histories keep their meaning.} Every saga's first record pins
the version of the state-machine rules it was admitted under (its
\emph{semantics version}, 1 originally and 4 today), stamped by the
engine rather than by the caller, and the fold branches on that version. A
change that makes a previously legal transition illegal is a new version;
sagas recorded under the old one go on folding exactly as they did. A version
the running build does not implement is refused, not guessed at. A replay
corpus of recorded histories, each with the projection it produced, runs on
every commit and fails if any history's meaning changes; at each of the three
version bumps so far, every earlier fixture was byte-identical. The corpus catches
a change of meaning only on the histories it contains, so a forgotten version bump
on an uncovered path would pass.
% source: pkg/saga/state.go (SemanticsVersion = 4); pkg/saga/testdata/fixtures (01-21 before the first bump, 22-26 added by the three bumps); make corpus

\subsection{Gates}
\label{sec:design-gates}
Gate requirements are resolved from a versioned policy when a saga is admitted
and recorded with it, so a saga is judged by the policy it was admitted under.
Each requirement is decided in one of two phases: \emph{before execution}, on
the facts the step proposes, or \emph{before release}, after it has run and
before its effect leaves the outbox, on the facts it was prepared with and the
results earlier steps published. A step \emph{declares} its facts in a prepare
record; the facts a pre-execution gate is asked about are its \emph{proposal};
each retry is a new \emph{attempt}. On the MCP edge, admission refuses a plan
whose irreversible step would be judged at a moment that can no longer stop it;
through the SDK, an irreversible step whose requirements are all release-phase
is admitted, and its body runs before they are decided (\cref{sec:limits}).
There are six gate types (schema, policy expression,
risk limit, frontier, validator and human), composed into one verdict.
% source: proto/janus/v1/common.proto (GateType)

Three properties make a gate more than a check the agent can talk its way past.
\emph{First, the fact set is closed}: a gate sees the values the step declared,
facts Janus derives from the record (\code{step.action},
\code{intent.principal}, \dots), and facts an \emph{earlier} step published under
\code{result.<step>.<key>}, a namespace no step can declare into. A payment
limit expressed against a balance reads a balance another step went and found;
the paying step never states it. \emph{Second, answers from outside are recorded
inputs}: a validator's opinion or a person's approval is a record, appended when
it arrives, and the verdict is re-decided over the log. The gate never calls
out, so a replay folds the approval that was given rather than asking again.
A human requirement is answered by a person: an answer from a participant is
refused whatever role it claims, except the system's refusal once the gate's
deadline has passed. It can demand an identity established through OIDC and a
WebAuthn step-up, and enforces separation of duty against the saga's principal
and its declared originator: declared by the participant that signed the begin,
not authenticated, so the check stops an honest mistake and not a liar. None of
these was exercised in the runs we report. A participant signs what
it asks to have recorded (an answer, a declaration, a result, a saga begin) with
an Ed25519 key its registered manifest declares; the daemon checks the signature
against the manifest version the saga pinned (the active one, for an answerer the
plan does not name), and answers and results keep it, so
the audit re-checks who answered from the log alone. A person's answer is signed
by the participant relaying it, and never by an agent when the daemon names its
human relayers or requires signatures. Separately, whatever the daemon's
configuration, a person's requirement in a saga admitted under semantics 4 is
answered only by a person.
% source: pkg/callersig; pkg/orchd/callers.go; pkg/gate/audit_signatures.go; pkg/saga/state.go (applyGateAnswer); pkg/gate/checks.go (checkHuman)
\emph{Third, an answer is bound to what it was asked about}: for a pre-execution
gate, the proposal an escalation was decided on (its facts and any sub-saga it
would delegate to) is pinned in the fold, every later verdict and the prepare on
that attempt must be about the same proposal, an answer counts only in the phase
its requirement is due and only once the question has been put, and the question
an answerer receives carries the facts (\cref{sec:findings}).

Each verdict is recorded with a \emph{decision provenance record} (DPR) naming
the requirements, the policy version and the grounds; the verdict cites the
record by event identifier, and the two share one durability barrier because no
decision stands between them.

\subsection{Effects and the interception edges}
\label{sec:design-edges}
An effect reaches the world through one of two paths, and they differ in how
much they enforce. On the \emph{MCP edge}~\cite{mcp}, a proxy in front of an
unmodified tool server turns each effectful call into a saga step and holds it
in the outbox; the tool server is never called on a refusal, and a held effect
is released only from a committed saga, under an idempotency key (I1, I4).
Delivery is at least once: the outbox records the key, but the MCP proxy does
not forward it, so a receiver that crashes mid-release may see a call twice.
The A2A middleware~\cite{a2a} checks the counterpart against the registry and
records both directions of a message, but does not gate or hold it.
% source: pkg/mcp/gateway.go (key not forwarded); pkg/outbox (at-least-once release); pkg/a2a/proxy.go (no saga)
On the \emph{SDK path}, a step's body is the effect and runs in the client: the
body runs only if the pre-execution gates pass, and the proposal, the verdict
and any refusal are durable before it does, but a client that ignored a
refusal would not be stopped: the log shows the refusal, and only the payment's
own record would show that the payment was made anyway. We call the first
\emph{mediation} and the second \emph{cooperative admission}, and
\cref{sec:eval-model} exercises the second.

An agent step can attach its own provenance: which model, how it was sampled,
the hash of what it output, the hashes of what it was shown, and the grounds it
gave. The daemon stamps the record's identity rather than trusting it, refuses a
decision kind that is the gate's to make, and writes it immediately before the
step's result, which cites it.

\subsection{The offline verifier}
\code{janus-verify} checks a log or an exported audit bundle against a trusted
root key: signatures, chain, Merkle roots, inclusion proofs, sequence
continuity, the writer-key chain and the envelope version, reporting each
finding by severity. \code{janus-gate audit}, given the same key, authenticates
the log first and then re-derives every recorded gate verdict from the inputs
recorded with it, reporting any the inputs do not support and how many answers
and results were recorded unsigned; without a key it says the log was not
authenticated. The verifier builds to identical bytes three
ways~\cite{reproduciblebuilds} and links only what it reads: 4.9\,MB for
darwin/arm64 and 5.2\,MB for linux/amd64 today. Linking a Protocol Buffers runtime took an
earlier, 7.0\,MB build to 16.3\,MB, so the payload it must parse (the
writer-key record) is CBOR, and a CI test fails if a denied package enters its
import graph.
% source: make repro (three builds identical); CGO_ENABLED=0 go build -trimpath -buildvcs=false: 4,904,242 bytes darwin/arm64, 5,160,278 linux/amd64 (2026-09-29); cmd/janus-verify import-graph test (TestTheVerifierLinksOnlyWhatItReads); 7.0 -> 16.3 MB measured when the protobuf runtime was linked
